% Transcription — fonds Alexandre Grothendieck, shelfmark 46, batch 3 (pages 41-60).
% Established from the facsimile archives/batches/46/batch-03.pdf (cut from a
% copy of 46.pdf fetched through the project relay,
% https://grothendieck-relay.onrender.com/source/46.pdf; PDF page k =
% archivists' page 40+k, no cover sheet), per the skill transcribe-grothendieck.
% The folder is ninety-four pages in five batches; this is the third.
%
% Pass: Opus 5.5 (claude-opus-5-5), 2026-09-26 — first pass, unchecked against the pages by a human.
%
% Pages skipped, and why:
%   - 41, 43, 45: three stencilled typed leaves, « -VI-6- », « -VI-5- »,
%     « -VI-4- », each headed « n°383 », scanned upside down: a duplicated
%     draft on graphs, forests and Coxeter transformations of a reflection
%     group (the Bourbaki-style « rédaction » numbering). They are the other
%     faces of the leaves of pages 42, 44, 46. Page 45 carries a small
%     interlinear correction and a marginal brace, of no identifiable hand
%     and no bearing on his mathematics. Unrelated versos; skipped.
%   - 47 (« - 3 - », « n° 399 », a corollary on divisible modules over a
%     principal ring) and 49 (« - 59 - », « n°357 », the normalised Haar
%     measure on Q_p): two further duplicated typed leaves of the same kind,
%     backs of pages 46 and 48, no ink of his. Skipped.
%   - 54: back of the cover of page 50; its only writing is an inscription
%     upside down at the foot, struck out, recorded in a \note after page 53.
% This is NOT the « Définition du morphisme Verschiebung » typescript of
% batch 2 (pages 23-26); no further leaf of that typescript occurs in pages
% 41-60.
% Covers, no \page marker, words as headings with a \note: 50 (« Progroupes
% et groupes algébriques quasi-compacts »), 55 (« Schémas de Brauer-Severi »).
%
% Pages summarised: none. The margins of pages 42, 44, 46, 48 and 52, written
% fast and on the slant, are partly compressed into \note{} with \ill{}.
%
% His pagination: none. Page 42 carries a circled « 6 » opening a numbered
% item (« Déf Système linéaire »); the items 1-5 are not in this batch
% (page 40, in batch 2, carries none). Page 56 is a plan numbered 1, 2, 3,
% 3 bis, 4.
%
% What the batch is about. Pages 42-48 continue the Brauer-Severi notes of
% page 40 (batch 2): corollaries — P a projective bundle iff it has local
% sections and an obstruction in H^2(S, O_S^*) vanishes; for S regular, iff
% locally trivial, iff its generic fibres are; local isotriviality; then
% (6) linear systems of divisors on X/S parametrised by a Brauer-Severi
% scheme P̌/S; the relation between morphisms X -> P̌ and linear systems;
% the morphism λ : Syst.Lin_{X/S} -> Pic_{X/S}, relatively representable by
% the « complete grassmannian » of a generalised Brauer-Severi scheme when
% X/S is proper with geometrically integral fibres; a class ξ of
% Pic_{X/S} represented by an invertible L on X iff P̌ is a projective
% bundle. Pages 51-53, under the cover « Progroupes et groupes algébriques
% quasi-compacts »: a diagram of functors Progr^!(C'_0) -> Gr(C) ≃ Gr(C')
% -> Gr qu.cpt(Sch/k), with the question whether it is fully faithful / an
% equivalence; every object of Pro Gr(C'_0) comes from a strict projective
% system, and the Hom of progroups is computed by lim_j colim_i Hom_gr;
% Lemme: a progroup whose transition maps are « ens. affines » is
% ens. affine. Pages 56-60, under the cover « Schémas de Brauer-Severi »
% (blue ink): a plan, « La théorie des Brauer-Severi-Châtelet globale »
% (definition, local isotriviality, projective bundles = Châtelet schemes
% with an L of rank 1 inducing O(1) on the fibres, dual bundle, the bundle
% of n-forms, openness of the Brauer-Severi locus, relation with Azumaya
% algebras and the global Brauer group); then the construction of the
% scheme of n-forms Phi^n(X) of a Brauer-Severi scheme X/Y, via the
% action of GP(r)_Y on P(S^n(O^{r+1})^v), its identification for
% X = P(E), and the Théorème: the sections of Phi^n(X) are the positive
% transversal divisors of degree n of X/Y; an example (divisorial ideals
% of the sheaf of total fractions R_0(X)) and a Proposition: for X flat over
% Y, X' ⊂ X closed is divisorial and Y-transversal at x iff X' is Y-flat at
% x and X' ⊗ k(y) is divisorial at x in X ⊗ k(y); the proof begins.
%
% Notation fixed once for the batch, following batch 2: the Brauer-Severi
% scheme and its dual are P, \check{P}; underlined sheaves keep their
% underline in math (\underline{L}, \underline{N}, \underline{E},
% \underline{\mathrm{Pic}}_{X/S}); structure sheaves \mathcal{O}. On
% pages 57-60 his letter for the scheme of n-forms (a looped capital, close
% to a Phi) is set \Phi^{n}(X), \Phi^{(n)}(X), and the divisors of degree n
% \mathcal{D}^{(n)}(X/Y). On pages 51-53 his categories are set
% \mathrm{Progr}^{!}(C'_0), \mathrm{Gr}(C), \mathrm{Gr\,qu.cpt}(\mathrm{Sch}_{/k}).
%
% Readings to check first: the margins of pages 42, 44, 46 and 52; the
% right-hand column of page 51; the second half of page 46; page 48.
% Apparatus: 179 \ill{}, 71 \uncertain{}.

\documentclass[11pt,a4paper]{article}
% Le préambule commun à toutes les transcriptions du fonds Grothendieck.
%
% Il tient en une idée : **l'incertitude se compose, elle ne se résout pas.**
% Une transcription qui lisse un mot douteux a détruit la seule chose pour
% laquelle on la faisait — un lecteur doit pouvoir savoir, à chaque endroit,
% ce qui était sur le feuillet et ce qui a été ajouté. Les six macros qui
% suivent sont donc l'appareil critique, et non de la décoration.
%
% Les mêmes commandes sont reconnues par `scripts/render.mjs`, qui produit la
% vue de lecture du site. Une macro ajoutée ici doit l'être aussi là-bas,
% faute de quoi le rendu échouera bruyamment — ce qui est voulu : un rendu
% silencieusement faux est pire qu'un rendu absent.

\ProvidesPackage{grothendieck}[2026/08/08 Transcriptions du fonds Grothendieck]

\usepackage{amsmath,amssymb,amsthm}
\usepackage{mathrsfs}
\usepackage{stmaryrd}
% Les diagrammes commutatifs sont partout dans ce fonds : c'est la langue
% même de la géométrie algébrique de Grothendieck, et une transcription qui
% les rendrait en prose ne transcrirait rien.
\usepackage{tikz-cd}
% Français par défaut — c'est la langue des feuillets — mais l'anglais est
% chargé aussi : la traduction et le résumé sont des documents anglais, et la
% typographie française (espaces avant les deux-points) y serait une faute.
% Ces documents basculent par \selectlanguage{english} en tête de corps.
\usepackage[english,french]{babel}
\usepackage{fontspec}
\usepackage{geometry}
\geometry{a4paper,margin=25mm,marginparwidth=28mm}
\usepackage{marginnote}
\usepackage{xcolor}
% ulem, not soul. `soul`'s \st and \ul are built on a character-by-character
% scan that XeLaTeX's font machinery breaks: the document fails with an
% undefined control sequence rather than degrading. ulem does the same two jobs
% and is Unicode-engine safe. `normalem` stops it hijacking \emph, which the
% transcriptions use for Grothendieck's own emphasis.
\usepackage[normalem]{ulem}
\usepackage{hyperref}
\hypersetup{colorlinks=true,linkcolor=black,urlcolor=black,pdfborder={0 0 0}}

% --- Métadonnées du lot -----------------------------------------------------
% Reprises telles quelles par le script de rendu, qui les affiche en tête de
% la vue de lecture. Elles fixent ce que le fichier prétend couvrir : sans
% elles, une transcription flotte sans point d'attache dans un fonds de seize
% mille feuillets.

\newcommand{\folder}[1]{\gdef\grothFolder{#1}}
\newcommand{\batch}[1]{\gdef\grothBatch{#1}}
\newcommand{\leaves}[2]{\gdef\grothFirst{#1}\gdef\grothLast{#2}}
\newcommand{\foldertitle}[1]{\gdef\grothFoldertitle{#1}}
\gdef\grothFolder{?}\gdef\grothBatch{?}\gdef\grothFirst{?}\gdef\grothLast{?}\gdef\grothFoldertitle{}

% --- Le résumé --------------------------------------------------------------
%
% Il ouvre la lecture modernisée, sous le titre « Résumé », et il est composé
% en petit. Ce n'est pas un ornement : c'est le seul passage du document qui
% ne soit pas des mathématiques, et un lecteur doit pouvoir le reconnaître
% avant de l'avoir lu — puis le sauter s'il connaît déjà le sujet, ou s'y
% arrêter s'il ne le connaît pas. Le filet qui le suit marque la reprise du
% corps.

\newenvironment{resume}
  {\small\setlength{\parskip}{0.45em}}
  {\par\medskip\hrule\medskip}

% --- Mots-clés ---------------------------------------------------------------
%
% En anglais, en clôture du résumé de la lecture modernisée : le vocabulaire
% moderne sous lequel chercher ce que les feuillets construisent. Le manifeste
% du site (`npm run manifest`) les extrait pour étiqueter la cote entière —
% cette ligne est la source unique des tags, il n'y a pas de fichier à côté.
\newcommand{\keywords}[1]{\par\smallskip\noindent
  {\footnotesize\textsc{Keywords} --- \textit{#1}}\par}

% --- L'appareil critique ----------------------------------------------------

% Le numéro de feuillet, en marge. C'est l'ancre qui permet de régler un
% désaccord sur une formule en regardant *un* feuillet plutôt qu'une chemise
% de six cents. Le site s'en sert aussi pour tourner les pages du fac-similé
% quand on fait défiler la transcription.
\newcommand{\leaf}[1]{%
  \marginnote{\footnotesize\color{gray!70}\textsf{#1}}%
  \label{leaf:#1}\ignorespaces}

% Illisible : on ne devine pas. Un mot inventé qui se lit comme les autres est
% le pire résultat possible.
\newcommand{\ill}{\textcolor{red!60!black}{[\,\ldots\,]}}

% Lecture douteuse : le mot est donné, et signalé comme incertain.
\newcommand{\uncertain}[1]{\uline{#1}}

% Ajout de l'éditeur — une lettre manquante, un mot sous-entendu.
\newcommand{\add}[1]{\textcolor{blue!50!black}{[#1]}}

% Biffé par Grothendieck lui-même. Ce qu'il a rayé fait partie du document :
% une rature dit souvent où la pensée a changé de direction.
\newcommand{\struck}[1]{\sout{#1}}

% Note du transcripteur, sur l'état matériel du feuillet.
\newcommand{\note}[1]{\par\smallskip{\small\color{gray!60!black}\textit{[#1]}}\par\smallskip}

% Note marginale de Grothendieck, distincte du corps du texte.
\newcommand{\marginal}[1]{\par\smallskip\noindent\hspace{1em}%
  {\small\color{gray!60!black}#1}\par\smallskip}

% --- Titre ------------------------------------------------------------------

\AtBeginDocument{%
  \begin{center}
    {\large\bfseries Fonds Alexandre Grothendieck --- cote n\textsuperscript{o}~\grothFolder}\\[2pt]
    {\itshape\grothFoldertitle}\\[4pt]
    {\small feuillets \grothFirst--\grothLast\ (lot \grothBatch)}\\[2pt]
    {\footnotesize\color{gray!60!black}Transcription établie d'après le fac-similé numérisé
     par l'Université de Montpellier.\\
     Les lectures incertaines sont soulignées, les illisibles notées [\,\ldots\,],
     les ajouts entre crochets.}
  \end{center}
  \vspace{1em}}

\endinput


\folder{46}
\batch{3}
\pages{41}{60}
\dating{[vers 1966-vers 1972]}
\watermark{Édition de démonstration}
\foldertitle{Schémas en groupes et fibrés principau[x] (Divers) : notes manuscrites (s.d.), tapuscrit (s.d.)}

\begin{document}

\page{42}

\emph{Notion loc.\ triviale :}

[\emph{Corollaire} \quad $P$ fibré projectif $\Longleftrightarrow$ $P$ admet
loc.\ une section, (i.e.\ est loc.\ \add{triviale} \struck{(fibré projectif)}),
et une certaine obstruction dans $H^2(S, \mathcal{O}_S^{*})$ est nulle\ldots]
\quad \uncertain{en remarque}.
\note{« triviale » est écrit au-dessus de « (fibré projectif) », biffé}

\emph{Corollaire} \quad Si $S$ régulier, $P$ fibré projectif ssi il est loc.\
trivial, \struck{\ill{}} ou encore ssi ses fibres \uncertain{génériques} le sont.

\quad Autre façon de procéder, valable si $S$ est strictement
\add{loc.}\ factoriel [i.e.\ tout \uncertain{schéma} étale dessus est
\add{loc.}\ factoriel] ----

\marginal{\uncertain{peut} \ill{} \uncertain{aussi} \ill{} \uncertain{$p^2$}}
\note{les deux corollaires et le paragraphe « Autre façon » sont embrassés à
gauche par une grande accolade, avec la note marginale écrite de biais,
presque verticalement ; lecture très incertaine. « loc. » est ajouté deux fois
au-dessus de la ligne}

\emph{Corollaire} \quad \struck{Loc} \uncertain{Loc.\ isotrivialité}.

(6)\note{le « 6 » est cerclé, dans la marge, sous un trait qui barre toute la
largeur de la page}
\quad \emph{Déf} \quad \emph{Système linéaire} \add{\struck{de} div.\ sur
$X/S$} (\emph{paramétré} par un schéma $\check{P}/S$ :
\uncertain{i.e.}\ loc.\ \uncertain{fppf} un syst.\ lin.\ de div.\ paramétré par un
fibré projectif. \struck{Système linéaire}
\quad [$X \to S$ \uncertain{f.\ plat} et loc.\ prés.\ finie]

\struck{Dualité que} \uncertain{Dualité} pour $X/S$ et un \ill{},
\uncertain{schéma} de Brauer-Severi : \emph{Système linéaire} (\uncertain{cas
complet}). Classe d'\uncertain{équivalence}\ldots
\note{la fin de la page est une suite de notes courtes, disposées sans ordre
évident ; l'ordre de lecture retenu ici est incertain}

\page{44}

a) \emph{Relation} \struck{\ill{}} Morphismes $X \to \check{P}$ et syst.\ lin.\ de
div.\ \ill{} \ill{} paramétrés par $\check{P}$.
\note{le « a) » n'est pas écrit ; le signe au-dessus du $\check{P}$ de la
première ligne est illisible}

b) \struck{Applications} Syst.\ lin.\ de div.\ « \uncertain{stricts} » \ill{} \ill{}
\ill{} paramétrés par $\check{P}$, et \uncertain{relations} \ill{} sur $S$.

$X \to \check{P}$ une petite propriété supplémentaire \emph{de nature locale}
[qui est inutile si $X$ est \uncertain{régulier}, \struck{\ill{}} ou strictement
loc.\ factoriel].

c)
\[
  \underline{\mathrm{Syst.Lin}}_{X/S} \xrightarrow{\ \lambda\ }
  \underline{\mathrm{Pic}}_{X/S}
\]
morphisme naturel (ayant d'ailleurs une section canonique). Si $X/S$ est
propre à fibres géom.\ intègres, alors $\lambda$ est rel.\ représentable comme
la grassmannienne complète d'un schéma de Brauer-Severi généralisé (attention,
il y a une question d'effectivité de données de descente sur un
\struck{fibré} fibré projectif relatif à un Module non loc.\ libre !].

Si, au lieu de fibres géom.\ intègres, on suppose seulement
$g_{*}(\mathcal{O}_X) \simeq \mathcal{O}_S$ \uncertain{univ.}, on trouve encore
un morphisme représentable par un \uncertain{ouvert} d'un fibré
\struck{projectif} de Brauer-Severi généralisé\ldots

\marginal{\uncertain{généralisation} \ldots\ $X \to S$ \uncertain{propre}
\ldots\ \ill{} \ill{} \ldots\ $\mathrm{Pic}(X/S)$ \ldots\ \uncertain{loc.}\
\ill{}}\note{cinq lignes serrées dans l'angle inférieur gauche, écrites de
biais, en partie barrées ; un mot souligné, peut-être un nom propre, ne se lit
pas}

\page{46}

\emph{Donc} \quad $\underline{\mathrm{Syst\,Lin\,div}}_{X/S}$ est représentable
ssi $\underline{\mathrm{Pic}}_{X/S}$ l'est.

\emph{Donnons-nous} un syst.\ linéaire paramétré par $\check{P}/S$ de Br.\ Sév.,
et considérons le
\[
  \underline{L} \qquad\qquad
  \underline{L} \otimes_{\mathcal{O}_S} \mathcal{O}_{\check{P}}(1) \otimes \underline{N}
\]

\begin{tikzcd}[column sep=small, row sep=small]
X \arrow[r, no head] \arrow[d, no head] & X \times_S \check{P} \arrow[d, no head] \\
S \arrow[r, no head] & \check{P}
\end{tikzcd}

section correspondante $\xi$ de $\underline{\mathrm{Pic}}_{X/S}$. Je dis que
$\xi$ \emph{est représentable par un $\underline{L}$ inversible sur $X$} ssi
[$f : X \to S$ propre, $f_{*}(\mathcal{O}_X) = \mathcal{O}_S$ \uncertain{univ.}]
$\check{P}$ \emph{est un fibré projectif} (\ill{} \ill{} par construction
\ill{})
\note{le premier $\underline{L}$ est au-dessus de $X$, le produit tensoriel
au-dessus de $X \times_S \check{P}$ ; les traits du carré n'ont pas de
pointes}

pour la condition est \uncertain{suffisante}, \ill{} \ill{}\ \struck{\ill{}} \ill{}
\ill{}, \ill{} \uncertain{facile} est nécessaire. \ill{} donc un Module de
prés.\ finie $Q$ sur $S$, et un \add{\uncertain{épi}}morphisme
$\check{P} \to \mathbb{P}(Q)$, je dis que c'est un isom.\ \ill{} \ill{}
$\mathbb{P}(E)$, où $E$ est un quotient loc.\ libre de $Q$. \ill{} \ill{} par
descente, type \ill{} \ill{} principe\ldots

\emph{Remarque} \quad Partant d'un Br.\ Sév.\ \ill{} \ill{}, on construit un
$\xi$ \ill{} \ill{} \ill{} \ill{} $\underline{L}$ \ill{} \ill{} \ill{}.

\marginal{\ill{} \ill{} $g_{*}(\mathcal{O}_X) \simeq \mathcal{O}_S$, de prés.\
finie, et $g$ \ill{}}\note{en haut de la marge gauche, de biais, en partie
barré}
\marginal{Si $X/S$ admet \ill{} \ill{} \ill{} \ill{} \ill{} \ill{} \ill{}
\ill{} $X$ \ldots}\note{en bas de la marge gauche, en cinq lignes obliques,
reliées par un trait vertical au paragraphe « Remarque »}

\page{48}

\[
  \underline{N} \simeq \underline{L} \otimes \mathcal{O}_{\check{P}}(1)
\]
\[
  \underline{N} \otimes_{\mathcal{O}_S} \underline{L}^{-1} = \mathcal{M}
\]
\struck{\ill{}} inversible sur $X \times_S \check{P}$, loc.\ au dessus de $S$
il provient d'un module inversible sur $\check{P}$, d'ailleurs de
\emph{façon canonique}, puisque $f_{*}(\mathcal{O}_X) = \mathcal{O}_S$ [pour
la \ill{} \ill{} que ce soit \uncertain{vrai} \uncertain{univ.}!], je \ill{}
\uncertain{supposer} \ill{} \ill{} \ill{}].
\note{la page est écrite par-dessus un texte dactylographié qui transparaît du
verso ; la fin, après « $f_{*}(\mathcal{O}_X) = \mathcal{O}_S$ », est serrée
et ne se lit qu'en partie}

\section*{Progroupes et groupes algébriques quasi-compacts}

\note{inscrit en haut à droite d'une chemise (page 50), qui ne reçoit pas de
numéro de page ; le premier mot de la deuxième ligne est surchargé, lecture
« \uncertain{groupes} » incertaine}

\page{51}

\begin{tikzcd}[column sep=small, row sep=small, nodes={font=\scriptsize}]
\mathrm{Progr}^{!}(C'_0) \arrow[r, "\alpha"] \arrow[d, hook, "i"'] & \mathrm{Gr}(C) \arrow[r, "\simeq", "\beta"'] & \mathrm{Gr}(C') \arrow[r, "\sigma"] \arrow[d, hook, "j"] & \mathrm{Gr\,qu.cpt}(\mathrm{Sch}_{/k}) \\
\mathrm{Progr}(C'_0) \arrow[rr, "k"'] & & \mathrm{Gr\,Pro}(C'_0) &
\end{tikzcd}

\note{au-dessus de $\beta$, « équiv. » ; au-dessus de $\sigma$, « pl.\ fid » ;
à côté des flèches $i$ et $j$, « pl.\ fid ». Les flèches $i$ et $j$ sont
dessinées comme des signes d'inclusion couchés. Après
$\mathrm{Gr\,qu.cpt}(\mathrm{Sch}_{/k})$, une flèche biffée avec son
étiquette. Dans l'angle supérieur gauche, encadré : « $! =$ ens.\ affine »,
et deux lignes biffées, « \ill{} syst.\ \uncertain{projectifs} \ill{} »}

$\beta$ équivalence car \struck{$\mathrm{Cart}(C) \simeq \mathrm{Cart}(C')$}
$C \cong C'$

$\sigma$ pl.\ fid trivial sur $\mathrm{Gr}(C) \to
\mathrm{Gr\,qu.cpt}(\mathrm{Sch}_{/k})$

$i$, $j$ pl.\ fid \struck{car} trivial [N.B.\ $C' \subset \mathrm{Pro}(C'_0)$
est pl.\ fid et commute aux $\varprojlim$ finies]

\[
  k \text{ pl.\ fid} \Longrightarrow \alpha \text{ pl.\ fid}
  \Longleftrightarrow
  \bigl[\sigma\alpha : \mathrm{Progr}^{!}(C'_0) \to
  \mathrm{Gr\,qu.cpt}(\mathrm{Sch}_{/k}) \text{ pl.\ fid.}\bigr]
\]
\note{le $k$ de tête est entouré, relié à « vrai », encadré}

$\sigma\alpha$ une équivalence $\Longleftrightarrow$ $\sigma$ et $\alpha$ des
équiv.\ $\Longrightarrow$ \struck{\ill{}} \ill{} schéma en groupes
qu.-compact sur $k$ admet \ill{} \ill{} \emph{affine}

\quad $\Downarrow$

$k$ une équivalence $\Longrightarrow$ $\alpha$ une équivalence
\note{sous « $k$ une équivalence », un « \uncertain{douteux} » entouré ; sous
la flèche, « pourvu que \struck{\ill{} \ill{}} \ill{} \ill{} tel que le
\uncertain{problème} \ill{} soit ens.\ affine, ainsi que \ill{} ens.\ affine »}

\note{à droite, sous l'énoncé « \ldots\ admet \ldots\ affine », et relié par
une longue ligne à « des équiv. » : « dans un schéma en groupes de t.f.\ sur
$k$ », précédé d'une double flèche vers la gauche et de « \ill{} \ill{} du
travail »}

\page{52}

Tout objet de $\mathrm{Pro\,Gr}(C'_0)$ provient d'un système \emph{projectif
strict}. Considérons deux tels $(G_i)_i$, $(H_j)_j$ et montrons que
\[
  \varphi : \varprojlim_j \varinjlim_i \mathrm{Hom}_{\mathrm{gr}}(G_i, H_j)
  \longrightarrow
  \Bigl(\varprojlim_j \varinjlim_i \mathrm{Hom}(G_i, H_j)\Bigr)_{\mathrm{mult.}}
\]
est \emph{bijectif}. En effet, comme pour $i$, $j$ fixés,
$\mathrm{Hom}_{\mathrm{gr}}(G_i, H_j) \to \mathrm{Hom}(G_i, H_j)$ est injectif, on
voit que $\varphi$ est \emph{injectif}. \struck{Reste à} Montrons que c'est
\emph{surjectif}, \struck{\ill{}} donc que si
$(u_j) \in \varprojlim_j \varinjlim_i \mathrm{Hom}(G_i, H_j)$ est multiplicatif,
alors les
\[
  u_j : (G_i)_i \longrightarrow H_j
\]
\struck{proviennent} \uncertain{donnent} des $u_j^{(i)} : G_i \to H_j$
\emph{mult.}
\struck{\ill{}} [si c'est vrai pour un, c'est vrai pour tous --].

\begin{tikzcd}[column sep=small, row sep=small]
G_j \times G_j \arrow[r] & G_i \times G_i \arrow[r] \arrow[d] & H \times H \arrow[d] \\
& G_i \arrow[r, no head] & H
\end{tikzcd}

\note{au-dessus de $G_j \times G_j$, un « $G_{j} \times G_{j}$ » biffé}
diagramme \uncertain{devient} comm.\ pour \uncertain{composé} avec
$G_j \times G_j \to G_i \times G_i$ pour $i$ grand, or
\struck{$G_j \times G_j \to G_i$} \ill{} est un \emph{épim.} de schémas
donc l'\uncertain{était} \uncertain{comme} \uncertain{avant}.

\marginal{N.B.\ \ill{} \ill{} \ill{} \ill{} \ill{} \ill{} \ill{}
\uncertain{transition} \ill{} \ill{} \ill{}}\note{dans l'angle inférieur
gauche, en six lignes obliques, dont une biffée ; souligné deux fois}

\page{53}

\emph{Lemme} \quad Tout progroupe \struck{ens.\ affine} $(G_i)_{i \in I}$ tel
que \ill{} \ill{} \ill{} soient ens.\ affines, est ens.\ affine.
\note{le lemme est marqué d'un trait vertical à gauche}

\struck{Dém.} On peut supposer $(G_i)_{i \in I}$ système projectif strict. Par
hyp.\ il existe un syst.\ projectif de schémas de prés.\ finie
\add{ens.\ affine} $(X_j)_{j \in J}$ qui lui est isomorphe comme proschéma. On
peut supposer les $X_j \to X_0$ affines.

On a

\begin{tikzcd}[column sep=small, row sep=small]
(G_i) \arrow[r] \arrow[d, no head] & G_{i_0} \arrow[d] \\
(X_j) \arrow[r, no head] & X_0
\end{tikzcd}

\note{la flèche de $G_{i_0}$ vers $X_0$ est oblique ; le trait vertical
entre $(G_i)$ et $(X_j)$ est très court et sans pointe}

Je dis que les $G_{i_1} \to G_{i_0}$ \add{$(i_1 \geq i_0)$} sont \emph{affines}.
En effet, \struck{$(G_i) \to$} $(X_j) \simeq (G_i) \to G_{i_1}$ se factorise par
$(X_j) \to X_{j_1} \to G_{i_1}$, rendant

\begin{tikzcd}[column sep=small, row sep=small]
(G_i) \arrow[r] \arrow[d, leftrightarrow] & G_{i_1} \arrow[r] & G_{i_0} \arrow[d] \\
X = (X_j) \arrow[r] & X_{j_1} \arrow[r] \arrow[u] & X_0
\end{tikzcd}

commutatif du carré $Q$, et
\note{la lettre $Q$ est écrite à l'intérieur du carré de droite}
\[
  G_{i_2} \longrightarrow X_{j_1} \longrightarrow G_{i_1} \longrightarrow G_{i_0}
  \qquad \Big/ \qquad
  X = \varprojlim G_i \longrightarrow G_{i_1} \longrightarrow G_{i_0}
\]
\note{sous $X_{j_1} \to G_{i_1}$, « [affine] » ; un arc de $X_{j_1}$ à
$G_{i_0}$ porte « affine » ; sous $G_{i_2} \to X_{j_1}$, « \uncertain{immersion}
\ill{} \ill{} \ill{} \ill{} » et un arc de $G_{i_2}$ marqué « fid.\ plat ».
Dans la seconde suite, « fid plat » au-dessus de $X \to G_{i_1}$ et un arc
« affine » sous $G_{i_1} \to G_{i_0}$. Dans la marge gauche, un long trait
courbe et un mot souligné, \ill{}}

\note{au dos de la chemise (page 54), seul, à l'envers au bas du feuillet et
barré de trois traits : « “Revêtements” étales de groupes
\uncertain{infinis} »}

\section*{Schémas de Brauer-Severi}

\note{inscrit en haut à droite d'une chemise bleue (page 55), qui ne reçoit
pas de numéro de page}

\page{56}

\emph{La théorie des Brauer-Severi-Châtelet globale.}

1. Définition d'un schéma de Brauer-Severi sur $Y$ : \struck{\ill{}}
\uncertain{loc.}\ $S$-isomorphe à $\mathbb{P}^r_Y$ \add{(i.e.\ \ill{} \ill{} \ill{}
$Y$ \ill{} Châtelet)}. [\ill{} \ill{} propre, plat sur $Y$, à fibres qui sont
des Brauer-Severi]. Possibilité de la $S$-descente, propriétés.
Loc.-isotrivialité.

2. Fibré projectif (associé à un faisceau loc.\ libre) $=$ schéma de Châtelet
sur $Y$ muni d'un $\underline{L}$ de rang $1$ qui sur chaque fibre induit
$\mathcal{O}(1)$ (\uncertain{bon} \uncertain{modulo}.) ;
\emph{Corollaire 1}, si $Y = \operatorname{Spec}(\mathcal{O})$, $\mathcal{O}$
local complet, alors Châtelet $=$ Brauer-Severi. \emph{Corollaire 2} Si $Y$
est \uncertain{régulier}, alors Br.\ Sév.\ loc.\ trivial $=$ \struck{fibré}
projectif $=$ Châtelet qui \struck{\ill{} \ill{}} \add{est loc.\ trivial} \ill{}
\ill{} \ill{}.
\note{au-dessus de « Corollaire 1 », ajouté : « \ill{} \ill{} Br.\ Sév.\ loc.\
triv. » ; la fin de l'énoncé du corollaire 2 est biffée d'un long trait, et
une ligne biffée suit : « \struck{Cor.\ 3 S'il existe un diviseur loc.\
principal \ill{}} »}

\marginal{\uncertain{Construction} des \uncertain{diviseurs}
« \uncertain{linéaires} »}

3. Fibré dual, (interprétation de ses sections comme diviseurs loc.\
principaux sur $X$ qui sont « \uncertain{hyperplans} »). \emph{Cor.~1} Si $X$
\uncertain{admet} une section, c'est un fibré projectif [$Y$ noeth.\ \ldots]
\quad $2$ façons de le prouver, \ldots\ $X$ \add{de Br.-Sév.}\ \ill{} une
section, \ill{} \ill{} \ill{} \ill{} donnée d'un \struck{fibré} faisceau loc.\
libre de rang $r$, et d'un sous faisceau loc.\ libre de rang $r-1$, loc.\
facteur direct\ldots \quad \emph{Cor.~2} \struck{Autre} Critère de loc.\
trivialité : existence de sections locales\ldots

\emph{3 bis} \quad Fibrés des $n$-\uncertain{formes} associé à
\struck{\ill{}} $X$, ses sections sur $Y$ sont les diviseurs loc.\ principaux
sur $X$ dont la trace sur toute fibre (ou une fibre, si $Y$ est connexe)
est de degré $n$.
\note{« 3 bis » est écrit en bas de la page, à l'encre plus foncée, après le
§~4 et l'accolade}

4. \struck{\ill{}} Soit $X$ propre et plat sur $Y$. \struck{\ill{}} Si
$f^{-1}(y)$ est un Br.-Sév., alors il existe un voisinage ouvert $U$ de $y$
t.q.\ $X|U$ soit Br.\ Sév.
\emph{Cor.} Châtelet $=$ Brauer-Severi.
\quad On prouve d'abord qu'il y a un $U$ t.q.\ $X|U$ soit Châtelet, puis on
\emph{considère} \ill{} \ill{} Châtelet $\Rightarrow$ Br.\ Sév., grâce à la
considération de $\underline{\mathrm{Isom}}_Y(\mathbb{P}^r_Y, X)$.

$\lbrace$Relations avec les faisceaux d'algèbres simples centrales\ldots\ Groupe de
Brauer global. [Consulter le cours de Serre\ldots]$\rbrace$

\page{57}

Soit $\underline{E}$ faisceau loc.\ libre de rang $r$. Alors
\struck{$GP(E)$ \ill{} opère sur} $\mathbb{P}(S^n(\underline{E}))$
\add{et $\mathbb{P}(S^n(\underline{E})^{\vee}) = \mathbb{P}(S^n(\check{\underline{E}}))$
\ill{} ne dépend que de $\mathbb{P}(\underline{E})$}
\struck{\ill{} \ill{}} et $GP(\underline{E})$ opère sur
$\mathbb{P}(S^n(\underline{E}))$ \add{\ill{}
$\mathbb{P}(\check{S}(\underline{E})^{\vee})$} [via
$GP(\underline{E}) \to GP(S^n_{\mathcal{O}_Y}(\underline{E}))$]. Faisant
$\underline{E} = \mathcal{O}_Y^{r+1}$, on trouve que $GP(r)_Y$ opère dans
$\mathbb{P}(S^n(\mathcal{O}_Y^{r+1})^{\vee})$ \add{via
$GP(r)_Y \to GP(S^n_{\mathcal{O}_Y}(\mathcal{O}_Y^{r+1}))$}, donc \ill{}
fibré principal homogène sous $GP(r)_Y$, \add{et \ill{} \ill{} \ill{}
schéma de Br.-Sév.\ $X$} \struck{\ill{} \ill{}} \ill{} fibré
\add{\uncertain{associé}} \struck{\ill{}} $S$-isomorphe : $\mathbb{P}(S^n(\mathcal{O}_Y^{r+1}))$,
qui est donc un schéma de Brauer-Severi. Il est loc.\ trivial (trivial) si $X$
\struck{\ill{} \ill{} \ill{} \ill{}} \add{l'est}. \add{On l'appelle le}
\emph{schéma des $n$-formes de} $X$, on le dénote par $\Phi^{n}(X)$.
Compatibilité avec images inverses.
\note{le début de la page est surchargé d'ajouts interlinéaires et de
passages biffés ; l'ordre de lecture retenu ici est incertain. Dans la
marge gauche, « \uncertain{n.} \ldots »}

Prenons d'abord $X = \mathbb{P}(\underline{E})$. Alors
\[
  \Phi^{n}(X) \overset{\text{can.}}{\simeq}
  \mathbb{P}\bigl(\underline{S}^{n}_{\mathcal{O}_X}(\check{\underline{E}})^{\vee}\bigr)
\]
[démonstration : \ill{} \ill{} \ill{} \ill{} des foncteurs, compatibles
\ill{} images inverses\ldots, qui \uncertain{admettent} une représentation, et
par suite une notion de fibré associé ; les autres \struck{fibrés} notions
\uncertain{s'en} \uncertain{déduisent} \ill{} \ill{} les foncteurs].

\struck{Sections de $\Phi$} Alors une section de $\Phi^{(n)}(X)$ est donc
définie par \struck{\ill{}} la donnée d'un \struck{\ill{}}
\add{\uncertain{homomorphisme}} \uncertain{surjectif} de
$\underline{S}^{n}(\check{\underline{E}})^{\vee}$ sur un \struck{\ill{}}
faisceau inversible, donc correspond à un hom.\ « \uncertain{ne s'annulant
nulle part} » d'un fibré inversible dans $\underline{S}^{n}(\underline{E})$,
i.e.\ une section de
\struck{\ill{}} $(\underline{S}^{n}(\underline{E}))^{*}/\mathcal{O}_Y^{*}$ ; or une
\note{sous $(\underline{S}^{n}(\underline{E}))^{*}$, une accolade : « germes
des sections de $\underline{S}^{n}(\underline{E})$ \uncertain{qui ne
s'annulent pas} »}

\page{58}

section $s$ de $\underline{S}^{n}(\underline{E})$ s'identifie à une section de
$\underline{L}_{\underline{E}}$ sur $X = \mathbb{P}(\underline{E})$, et
\struck{\ill{}} \add{définit} \struck{\ill{}} un \emph{diviseur}
\add{$\mathrm{div}(s)$} \struck{\ill{}} sur $X$ ; la condition que $s$
[sur $Y$] ne s'annule \emph{\uncertain{nulle part}}, signifie que ce diviseur
induit bien un diviseur sur chaque fibre, \struck{\ill{}} i.e.\ que $s$ ne
s'annule pas. D'ailleurs \struck{\ill{} \ill{}} $\mathrm{div}(s) =
\mathrm{div}(t)$, \struck{\ill{}} si et seulement si on peut trouver
$\rho \in \Gamma(X, \mathcal{O}_X^{*})$ t.q.\ $t = \rho s$, or $\rho$
s'identifie à un élément de $\Gamma(Y, \mathcal{O}_Y^{*})$,
(car $f_{*}(\mathcal{O}_X) = \mathcal{O}_Y$). Donc $\mathrm{div}(s) =
\mathrm{div}(t)$ $\Longleftrightarrow$ $s$ et $t$ définissent la même section
de $\underline{S}^{n}(\underline{E})^{*}/\mathcal{O}_Y^{*}$.

De plus, $s \mapsto \mathrm{div}(s)$ est compatible avec localisation sur $S$.
D'où : \emph{on a une application injective} de
\[
  \Phi^{(n)}(X)(Y) = \Gamma\bigl(Y, \underline{S}^{n}(\underline{E})^{*}/\mathcal{O}_Y^{*}\bigr)
\]
\emph{dans l'ens.\ des diviseurs} \emph{loc.\ principaux} \add{positifs}
\emph{sur $X$} \emph{qui induisent} \struck{un diviseur} \emph{sur chaque
fibre} \emph{un diviseur}, « \add{\uncertain{positif}} \emph{de degré $n$} ».
Appelons de tels diviseurs sur $X$ des « diviseurs transversaux de degré $n$ ».
Je dis que l'application précédente est \emph{surjective}, i.e.\ que tout
diviseur transversal $D$ de degré $n$ provient d'une section de
$\Phi^{(n)}(X)$. En effet, soit $\underline{L} = \underline{L}(D)$ le faisceau
inversible défini par

\page{59}

$D$, \struck{\ill{} \ill{}} la question étant locale \add{sur $Y$}, on peut
supposer que $\underline{L} \simeq \mathcal{O}_X(k)$ pour $k$ convenable,
d'autre part on voit aussitôt, puisque
$\underline{L}|f^{-1}(y) \simeq \mathcal{O}_{f^{-1}(y)}(k)$, que $k = n$. Donc
$D$ est le diviseur d'une section de $\mathcal{O}_X(k)$.

Ainsi, on a \struck{\ill{}} trouvé un isomorphisme \add{\uncertain{canonique}}
\[
  \Phi^{(n)}(\mathbb{P}(\underline{E})) \xrightarrow{\ \sim\ } \mathcal{D}^{(n)}(X/Y) .
\]
Cet isomorphisme est fonctoriel, compatible avec images inverses.
D'autre part, la notion de « diviseur \add{loc.\ principal positif}
transversal » est compatible à descente [car c'est la notion de : fibré
inversible, muni d'une section régulière qui est \uncertain{non} div.\ de $0$
et \uncertain{non} diviseur de $0$ sur chaque fibre]. De plus, la notion de
degré d'un diviseur sur un schéma algébrique de Severi-Brauer a un sens
\struck{[\ill{} \ill{} \ill{} \ill{}]}, invariante par extension du corps de
base, d'où la notion de diviseur \add{positif} \struck{\ill{}} transversal de
degré $n$ sur un sch.\ Brauer. Cette notion est compatible à la descente.
D'où le

\emph{Théorème} \quad Les sections de $\Phi^{(n)}(X)$ ($X$ un sch.\ Sév.\
Br.) sont les diviseurs \add{positifs} transversaux de degré $n$ de $X/Y$.

\page{60}

\emph{Exemple de diviseurs transversaux}. Soit $X$ \struck{de type fini}
sur $Y$, \struck{\ill{}} $R_0(X)$ le faisceau \uncertain{total} des fractions
de $\mathcal{O}_X$, \struck{$R_0$} Alors
$\Gamma\bigl(R_0(X)^{*}/\mathcal{O}_X^{*}\bigr) \simeq$ \struck{\ill{}}
idéaux fractionnaires divisoriels loc.\ principaux de $X$ $\simeq$ diviseurs
loc.\ principaux sur $X$.
\note{« idéaux fractionnaires » remplace, au-dessus de la ligne, un
« \struck{groupe} » et un « \struck{\ill{} des diviseurs} » biffés ; « loc. »
est ajouté au-dessus de « principaux »}
Soit $\underline{I}$ un idéal de $R_0(X)$, on donne des conditions pour qu'il
soit \struck{principal} \emph{divisoriel loc.\ principal transversal}
\uncertain{$\xi$}. On suppose $\underline{I}$ \uncertain{positif}, i.e.\
$\underline{I} \subset \mathcal{O}_X$, et on suppose \ill{} \emph{\ill{}} que
$\underline{I}$ est un idéal de $\mathcal{O}_X$, \uncertain{$\xi$} donc
correspond à un sous-schéma $X'$ de $X$.

\emph{Proposition} \quad Supposons $X$ plat sur $Y$, $X$, $Y$ loc.\ noeth.,
soit $X'$ un sous-schéma fermé de $X$. On suppose $X$ \emph{plat} sur $Y$.
Soit $x \in X'$, [$y$ son image dans $Y$]. Pour que $X'$ soit divisoriel et
transversal \struck{\ill{}} \add{$Y$-transversal} en $x$, il faut et il suffit
que l'on ait les conditions suivantes
\begin{quote}
(i) $X'$ est $Y$-plat en $x$

(ii) $X' \otimes k(y)$ est divisoriel \struck{\ill{}} en $x$ dans
$X \otimes k(y)$.
\end{quote}
\note{l'énoncé est marqué d'un trait vertical à gauche}

\emph{Démonstration}. \emph{\uncertain{Nécessité}} On peut supposer
$\mathcal{O}_{X'} = \mathcal{O}_X/\varphi\mathcal{O}_X$,
$\varphi \in \Gamma(X, \mathcal{O}_X)$, $\varphi$ \struck{\ill{} \ill{}
\ill{}} \add{une fonction} \struck{\ill{}} \emph{divisorielle}, et $\varphi$
induisant dans les fibres des fonctions divisorielles. Alors \ill{} \ill{}
\ill{} exacte
\[
  0 \to \mathcal{O}_X \xrightarrow{\ \varphi\ } \mathcal{O}_X \to
  \mathcal{O}_{X'} \to 0 ,
\]
\add{et si $z \in Y$} \struck{\ill{} \ill{}}
\[
  0 \to \mathcal{O}_X \otimes k(z) \xrightarrow{\ \varphi\ }
  \mathcal{O}_X \otimes k(z) \to \mathcal{O}_{X'} \otimes k(z) \to
\]
\note{au début de la dernière ligne, « \struck{$\to \mathcal{O}_{X'} \to$
\ill{}} » biffé ; la page s'arrête ici}

\end{document}
